% ============================================================
%  CAPÍTULO 6: RESULTADOS EMPÍRICOS
% ============================================================
\chapter{Resultados empíricos}
\label{cap:resultados}

% ------------------------------------------------------------
\section{Introducción}
% ------------------------------------------------------------

Este capítulo presenta los resultados de las cinco especificaciones estimadas:
modelo base con un retardo de espera (M1, efectos fijos y aleatorios), modelo
con dos retardos de espera (M2), modelo con persistencia del gasto (M3), y dos
modelos de umbral que reemplazan la variable de espera continua por una dummy
de espera superior a 100 días (M1-U y M3-U). Se evalúa en qué medida los
resultados apoyan las hipótesis de investigación formuladas en el
Capítulo~\ref{cap:introduccion}.

% ------------------------------------------------------------
\section{Resultados principales: Tabla de regresiones}
% ------------------------------------------------------------

La Tabla~\ref{tab:resultados} presenta los coeficientes estimados, errores
estándar robustos agrupados por CCAA, y los indicadores de ajuste para los
seis modelos estimados.

\begin{table}[htbp]
  \centering
  \caption{Resultados de los modelos de regresión de panel
           (Variable dependiente: $\log G_{it}$)}
  \label{tab:resultados}
  \begin{threeparttable}
    \scriptsize
    \setlength{\tabcolsep}{4pt}
    \resizebox{\textwidth}{!}{%
    \begin{tabular}{lS[table-format=-2.4]S[table-format=-2.4]
                       S[table-format=-2.4]S[table-format=-2.4]
                       S[table-format=-2.4]S[table-format=-2.4]}
      \toprule
      & {M1-FE} & {M1-RE$^*$} & {M2-FE} & {M3-FE} & {M1-U} & {M3-U} \\
      \midrule
      Constante &
        -5.1640* & -9.2143*** & -4.9325* & -5.1925* & -6.0774** & -5.7288* \\
        & {(2.8807)} & {(2.1142)} & {(2.7984)} & {(2.8603)} & {(2.9046)} & {(2.9345)} \\
      \addlinespace
      $\log\text{PIB\,pc}_{it}$ &
        1.0151*** & 1.5414*** & 0.9656*** & 0.9246** & 1.1367*** & 1.0014*** \\
        & {(0.3681)} & {(0.2298)} & {(0.3531)} & {(0.3556)} & {(0.3682)} & {(0.3606)} \\
      \addlinespace
      $\text{Espera}_{i,t-1}$ &
        0.0002 & 0.0005 & -0.0001 & -0.0000 & {---} & {---} \\
        & {(0.0006)} & {(0.0005)} & {(0.0006)} & {(0.0006)} & & \\
      \addlinespace
      $\text{Espera}_{i,t-2}$ &
        {---} & {---} & 0.0003 & {---} & {---} & {---} \\
        & & & {(0.0005)} & & & \\
      \addlinespace
      $\mathbb{1}(\text{Espera}_{t-1}>100)$ &
        {---} & {---} & {---} & {---} & 0.0926** & 0.0558** \\
        & & & & & {(0.0380)} & {(0.0224)} \\
      \addlinespace
      $\log G_{i,t-1}$ &
        {---} & {---} & {---} & 0.2104* & {---} & 0.2010* \\
        & & & & {(0.1168)} & & {(0.1166)} \\
      \addlinespace
      $D_{\text{COVID},t}$ &
        0.0906 & 0.1657*** & 0.0888 & 0.0780 & 0.1063* & 0.0880 \\
        & {(0.0589)} & {(0.0511)} & {(0.0571)} & {(0.0601)} & {(0.0582)} & {(0.0597)} \\
      \addlinespace
      $\text{Pob65}_{it}$ &
        0.0045 & -0.0608*** & 0.0176 & -0.0018 & -0.0114 & -0.0118 \\
        & {(0.0539)} & {(0.0199)} & {(0.0544)} & {(0.0495)} & {(0.0503)} & {(0.0461)} \\
      \midrule
      Observaciones & 118 & 118 & 101 & 116 & 118 & 116 \\
      $R^2$ & 0.2816 & 0.5484 & 0.2002 & 0.5841 & 0.3895 & 0.6271 \\
      \bottomrule
    \end{tabular}}
    \begin{tablenotes}[flushleft]
      \footnotesize
      \item Nota: Errores estándar robustos agrupados por CCAA entre paréntesis.
      \item $R^2$ corresponde al $R^2$ global (\textit{overall}) de cada modelo.
      \item $^*$ M1-RE presentado para comparación; test de Hausman ($\chi^2(4)=11{,}44$, $p=0{,}022$) prefiere FE.
      \item El número de observaciones varía según disponibilidad de variables
        (17 CCAA, 2016--2024, con datos faltantes).
      \item Significatividad: $*$ $p < 0{,}10$, $**$ $p < 0{,}05$,
        $***$ $p < 0{,}01$.
    \end{tablenotes}
  \end{threeparttable}
\end{table}

% ------------------------------------------------------------
\section{Test de Hausman: Elección entre efectos fijos y aleatorios}
\label{sec:hausman_resultados}
% ------------------------------------------------------------

El test de Hausman contrasta si los efectos individuales están correlacionados
con los regresores. El estadístico obtenido es $\chi^2(4) = 11{,}44$ con
$p$-valor de $0{,}022$. \textbf{Se rechaza la hipótesis nula} al 5\%, lo que
indica que los efectos individuales están correlacionados con los regresores y
el estimador de efectos aleatorios no es consistente. Por tanto,
\textbf{M1-FE es el modelo preferido} para la interpretación principal. No
obstante, los resultados de efectos aleatorios se presentan como referencia
para la comparación con la variación \textit{between}.

% ------------------------------------------------------------
\section{Contraste de H1: Elasticidad renta}
\label{sec:h1}
% ------------------------------------------------------------

\textbf{H1}: \textit{El gasto en seguros sanitarios privados es un bien de lujo
($\varepsilon_Y > 1$)}.

El coeficiente de $\log\text{PIB\,pc}$ en M1-FE es $\hat{\beta}_1 = 1{,}02$
(e.e.\ 0{,}37), significativo al 1\%. Un incremento del 1\% en la renta per
cápita se asocia con un incremento del \textbf{1{,}02\%} en el gasto en seguros
privados. El coeficiente es prácticamente igual a la unidad, lo que sitúa el
seguro sanitario privado en la frontera entre bien normal y bien de lujo. En el
modelo de efectos aleatorios la elasticidad es más elevada
($\hat{\beta}_1 = 1{,}54$, e.e.\ 0{,}23), pero el test de Hausman descarta la
consistencia de este estimador.

\textbf{Conclusión:} H1 queda \textbf{confirmada}. La elasticidad renta es
significativamente positiva y cercana o superior a la unidad, confirmando que
el seguro sanitario privado se comporta como un \textbf{bien de lujo} a escala
regional.

% ------------------------------------------------------------
\section{Contraste de H2: Efecto lineal de las listas de espera}
\label{sec:h2}
% ------------------------------------------------------------

\textbf{H2}: \textit{Tiempos de espera más largos inducen mayor gasto en
seguros privados}.

En M1-FE, el coeficiente de $\text{Espera}_{i,t-1}$ es $\hat{\gamma}_1 = 0{,}0002$
(e.e.\ 0{,}0006), con $p$-valor de $0{,}77$. El coeficiente tiene signo positivo
pero \textbf{no es estadísticamente significativo}. Bajo efectos aleatorios el
resultado es similar ($\hat{\gamma}_1 = 0{,}0005$, $p = 0{,}34$).

Este resultado sugiere que el efecto de las listas de espera sobre el gasto
\textbf{no opera de forma lineal} a través de la variable continua de días de
espera. Sin embargo, como veremos en la sección~\ref{sec:h4}, el efecto sí es
significativo cuando se especifica como variable umbral.

\textbf{Conclusión:} H2 \textbf{no queda confirmada} en su forma lineal.

% ------------------------------------------------------------
\section{Contraste de H3: Efecto diferido (dos retardos)}
\label{sec:h3}
% ------------------------------------------------------------

\textbf{H3}: \textit{El efecto de las listas de espera opera con estructura de
retardos temporal.}

En M2-FE, los coeficientes de los dos retardos son: $\hat{\gamma}_1 = -0{,}0001$
(e.e.\ 0{,}0006) y $\hat{\gamma}_2 = 0{,}0003$ (e.e.\ 0{,}0005). El test de
Wald para la hipótesis conjunta $H_0: \gamma_1 = \gamma_2 = 0$ produce:

\begin{equation}
  \chi^2(2) = 0{,}39 \quad (p = 0{,}82)
  \label{eq:wald_lags}
\end{equation}

\textbf{No se rechaza la hipótesis nula}. Los retardos de las listas de espera
no son conjuntamente significativos, lo que sugiere que el efecto lineal de los
días de espera no se acumula de forma diferida.

\textbf{Conclusión:} H3 \textbf{no queda confirmada}. El efecto diferido no es
significativo con la especificación lineal de retardos.

% ------------------------------------------------------------
\section{Contraste de H4: Efecto umbral}
\label{sec:h4}
% ------------------------------------------------------------

\textbf{H4}: \textit{Existe un efecto umbral tal que, por encima de 100 días de
espera, el impacto sobre el gasto privado es significativamente mayor.}

En M1-U, el coeficiente de $\mathbb{1}(\text{Espera}_{t-1}>100)$ es
$\hat{\delta} = 0{,}0926$ (e.e.\ 0{,}0380), con $t = 2{,}43$ y $p$-valor de
$0{,}017$. Este resultado implica que, cuando la espera del año anterior supera
los 100 días, el gasto en seguros privados es aproximadamente un
\textbf{9{,}3\% mayor}, siendo este efecto \textbf{estadísticamente
significativo al 5\%}.

El modelo M3-U, que incorpora además el logaritmo del gasto anterior, confirma
el efecto umbral ($\hat{\delta} = 0{,}0558$, e.e.\ 0{,}0224, $p = 0{,}015$),
indicando un incremento del 5{,}6\% en el gasto cuando la espera retardada
supera el umbral. La persistencia del gasto anterior también es significativa
al 10\% ($\hat{\rho} = 0{,}201$, $p = 0{,}088$).

Este hallazgo es central para la investigación: \textbf{el efecto de las listas
de espera no es lineal}, sino que opera como un mecanismo de ``disparo'' cuando
se supera un umbral crítico. Esto es coherente con el efecto \textit{exit} de
Hirschman: los hogares no ajustan gradualmente su comportamiento ante pequeños
incrementos de espera, sino que ``saltan'' al sector privado cuando perciben
que el sistema público ha cruzado un límite de tolerancia.

\textbf{Conclusión:} H4 queda \textbf{confirmada}. Existe un efecto umbral
significativo en los 100 días de espera.

% ------------------------------------------------------------
\section{Contraste de H5: Persistencia del gasto}
\label{sec:h5}
% ------------------------------------------------------------

\textbf{H5}: \textit{El gasto en seguros privados presenta inercia temporal
(el gasto del año anterior predice el gasto actual).}

En M3-FE, el coeficiente del logaritmo del gasto en seguros del año anterior es
$\hat{\rho} = 0{,}210$ (e.e.\ 0{,}117), con $t = 1{,}80$ y $p$-valor de
$0{,}075$. Dado que tanto la variable dependiente como el regresor están en
logaritmos, este coeficiente se interpreta como una \textbf{elasticidad}: un
aumento del 1\% en el gasto del año anterior se asocia con un incremento del
0{,}21\% en el gasto actual. En M3-U, el coeficiente es $\hat{\rho} = 0{,}201$
($p = 0{,}088$), confirmando el orden de magnitud.

Este resultado indica que existe \textbf{inercia moderada} en el gasto en
seguros privados: las CCAA que gastaron más el año anterior tienden a gastar
más en el período actual, incluso controlando por renta y listas de espera.
Esta persistencia puede reflejar: (i) renovación automática de pólizas,
(ii) formación de hábitos de consumo sanitario privado, o (iii) procesos de
difusión cultural del aseguramiento privado.

\textbf{Conclusión:} H5 queda \textbf{parcialmente confirmada} en el modelo de
efectos fijos con retardo (M3-FE). La elasticidad de persistencia de 0{,}21
es significativa al 10\% ($p = 0{,}075$), lo que sugiere inercia temporal
moderada en el gasto en seguros privados, aunque la evidencia no es concluyente
al nivel convencional del 5\%.

% ------------------------------------------------------------
\section{Impacto del COVID-19}
\label{sec:covid}
% ------------------------------------------------------------

El coeficiente de $D_{\text{COVID},t}$ es positivo en todos los modelos. En
M1-FE (modelo preferido), $\hat{\beta}_4 = 0{,}091$ (e.e.\ 0{,}059), positivo
pero no significativo al 5\% ($p = 0{,}13$). En el modelo de efectos aleatorios
el efecto es mayor y significativo ($\hat{\beta}_4 = 0{,}166$, e.e.\ 0{,}051,
$p = 0{,}002$), implicando un incremento del \textbf{16{,}6\%}. En M1-U, el
coeficiente es marginalmente significativo ($\hat{\beta}_4 = 0{,}106$,
$p = 0{,}071$).

Estos valores son exactamente los reportados en la Tabla~\ref{tab:resultados}:
M1-FE = 0{,}0906 (n.s.) y M1-RE = 0{,}1657$^{***}$.

Este efecto positivo, aparentemente contraintuitivo, es coherente con el
colapso percibido de la sanidad pública durante la pandemia y la búsqueda de
alternativas privadas. La significatividad del efecto depende de la
especificación utilizada.

% ------------------------------------------------------------
\section{Resumen de contraste de hipótesis}
\label{sec:resumen_hipotesis}
% ------------------------------------------------------------

La Tabla~\ref{tab:hipotesis} sintetiza el contraste de las cinco hipótesis.

\begin{table}[htbp]
  \centering
  \caption{Síntesis del contraste de hipótesis}
  \label{tab:hipotesis}
  \begin{threeparttable}
    \small
    \resizebox{\textwidth}{!}{%
    \begin{tabular}{clp{4cm}cc}
      \toprule
      \textbf{Hip.} & \textbf{Enunciado} & \textbf{Evidencia}
        & \textbf{Resultado} & \textbf{Modelo} \\
      \midrule
      H1 & Elasticidad renta $>1$ &
        $\hat{\beta}=1{,}02^{***}$; $p=0{,}007$ &
        \textbf{Confirmada} & M1-FE \\
      H2 & Espera $\uparrow \Rightarrow$ seguros $\uparrow$ (lineal) &
        $\hat{\gamma}=0{,}0002$ (n.s.) &
        No confirmada & M1-FE \\
      H3 & Efecto diferido (retardos) &
        Wald $\chi^2(2)=0{,}39$ ($p=0{,}82$) &
        No confirmada & M2-FE \\
      H4 & Efecto umbral (100 días) &
        $\hat{\delta}=0{,}093^{**}$; $p=0{,}017$ &
        \textbf{Confirmada} & M1-U \\
      H5 & Inercia del gasto &
        $\hat{\rho}=0{,}210^{*}$; $p=0{,}075$ &
        Parcialmente confirmada & M3-FE \\
      \bottomrule
    \end{tabular}}
    \begin{tablenotes}[flushleft]
      \footnotesize
      \item Errores estándar agrupados por CCAA.
      \item Test de Hausman: $\chi^2(4)=11{,}44$, $p=0{,}022$ $\Rightarrow$ FE preferido.
      \item Significatividad: $*$ $p < 0{,}10$; $**$ $p < 0{,}05$; $***$ $p < 0{,}01$.
    \end{tablenotes}
  \end{threeparttable}
\end{table}
